\section{Nguyên lý hoạt động}

\subsection{Quá trình ghi (Master Transmit)}
Viết START = 1 trong \texttt{i\_config\_reg} để bắt đầu giao dịch.

\texttt{I2C\_BUSY} được set lên trước khi START bị auto-clear.

Trình tự:
\begin{itemize}[leftmargin=1.2cm]
    \item START: SDA kéo xuống thấp khi SCL đang ở mức HIGH.
    \item Byte địa chỉ: 7-bit địa chỉ + 1-bit R/W = 0 (write). Master chỉ được đổi SDA khi SCL = LOW.
    \item Bit ACK (bit thứ 9): master nhả SDA (high-Z) và lấy mẫu ở pha SCL HIGH; slave kéo thấp $\Rightarrow$ ACK, nhả cao $\Rightarrow$ NACK.
    \item Các byte dữ liệu: mỗi byte gồm 8 bit MSB-first + 1 bit ACK tương tự.
    \item STOP: SDA lên HIGH khi SCL đang HIGH để nhả bus.
\end{itemize}

\texttt{TX\_DONE} được set sau khi byte cuối được gửi.

\texttt{I2C\_BUSY} chỉ clear khi STOP hoàn tất.

\textbf{Trường hợp lỗi/NACK:} nếu NACK xảy ra ở địa chỉ hoặc ở một byte dữ liệu, Master FSM thường set \texttt{TX\_ERR} và kết thúc giao dịch sớm (phát STOP hoặc ABORT tùy cấu hình), đồng thời báo ngắt để CPU xử lý retry hoặc báo lỗi thiết bị.

\subsection{Quá trình đọc (Master Receive)}
Tương tự truyền, nhưng bit R/W = 1.

Byte dữ liệu từ slave được lấy vào \texttt{o\_receive\_data}, với \texttt{o\_received\_data\_valid} = 1.

Khi byte cuối đến, \texttt{RX\_DONE} được set.

Ở chế độ đọc, master có thể trả ACK hoặc NACK cho từng byte dữ liệu từ slave (thường NACK byte cuối trước STOP).

Trong pha đọc, \textbf{nguồn điều khiển SDA đảo vai}: slave lái dữ liệu từng bit, còn master chỉ điều khiển bit ACK/NACK sau mỗi byte bằng cách kéo SDA thấp (ACK) hoặc nhả SDA (NACK) ở bit thứ 9.

\subsection{Repeated START}
Mục đích: đổi hướng hoặc truy cập thanh ghi nội bộ slave (ví dụ ghi địa chỉ thanh ghi rồi Sr rồi đọc dữ liệu) \emph{không} phát STOP giữa hai phase.

Trong lần START thứ hai, phải giữ nguyên địa chỉ slave theo đặc tả nhóm.

Trên bus, Sr là điều kiện START khi bus đang bận (không có STOP trước đó): SDA chuyển từ cao xuống thấp khi SCL = 1 sau phase trước.

Kịch bản kiểm thử điển hình: START $\rightarrow$ addr+W $\rightarrow$ ACK $\rightarrow$ (data) $\rightarrow$ Sr $\rightarrow$ addr+R $\rightarrow$ \ldots\ — không có STOP giữa hai START.

Điểm kỹ thuật cần kiểm: \textbf{I2C\_BUSY không được clear} giữa hai START; Bus FSM phải tạo START thứ hai trong khi vẫn đang chiếm bus (SCL/SDA không trở về trạng thái idle kéo dài).

